% !TeX root = ../../Book3.tex
\section{Dirichlet 卷积}
\label{b3:sec:2502}

“对所有约数求和”是一种线性变换。为了反演它，把一般两个算术函数按因数对求和，便得到与整数乘法适配的卷积。

\subsection{定义}
\label{b3:subsec:250201}

\begin{definition}\label{b3:def:25-convolution}
算术函数 \(a,b\) 的\term[dirichlet-juanji]{Dirichlet 卷积}[Dirichlet convolution]为
\((a*b)(n)=\sum_{d\mid n}a(d)b(n/d)\)。
\end{definition}
每个系数只有有限和，无需收敛假设。把它写成 \(\sum_{uv=n}a(u)b(v)\)，立即得到交换律。结合律来自两种括号都等于 \(\sum_{uvw=n}a(u)b(v)c(w)\)；对加法的分配律逐项成立。

\subsection{单位元}
\label{b3:subsec:250202}

函数 \(\varepsilon(1)=1\)、\(\varepsilon(n)=0\ (n>1)\) 为卷积单位元；常函数 \(\mathbf1(n)=1\) 则不是单位，而是约数求和算子：\((\mathbf1*a)(n)=\sum_{d\mid n}a(d)\)。这与数值逐点乘法的单位不同。

\ProseParagraphBreak

两个可乘函数的卷积仍可乘。若 \(m,n\) 互素，每个 \(mn\) 的约数唯一写作 \(d_1d_2\)，\(d_1\mid m,d_2\mid n\)。将两个函数的乘法公式代入有限和，便分成 \((a*b)(m)(a*b)(n)\)。互素假设在约数分离这一步使用。

\subsection{逆元}
\label{b3:subsec:250203}

\begin{proposition}\label{b3:prop:25-inverse}
算术函数 \(a\) 有卷积逆元，当且仅当 \(a(1)\ne0\)，逆元唯一。
\end{proposition}
\begin{proof}
若 \(a*b=\varepsilon\)，在一处得 \(a(1)b(1)=1\)，故必要。反之置 \(b(1)=1/a(1)\)，对 \(n>1\) 递推定义
\[
b(n)=-\frac1{a(1)}\sum_{d\mid n,\ d>1}a(d)b(n/d).
\]
和中 \(n/d<n\)，所以按整数大小递推合法，并逐个满足卷积方程。该公式也强制任何逆元具有相同值，故唯一。
\end{proof}
若 \(a\) 可乘，其逆元也可乘。可先分别在每个素数幂上递推得到逆值，再按互素分解乘起来；局部有限卷积乘积显示得到的函数与 \(a\) 卷积为 \(\varepsilon\)，由唯一性即为逆元。这里不必借用解析级数的倒数。

\subsection{Möbius 反演}
\label{b3:subsec:250204}

\begin{theorem}[Möbius 反演]\label{b3:thm:25-mobius-inversion}
对算术函数 \(a,b\)，等式 \(b(n)=\sum_{d\mid n}a(d)\) 等价于
\(a(n)=\sum_{d\mid n}\mu(d)b(n/d)\)。
\end{theorem}
\begin{proof}
\cref{b3:prop:25-mobius-sum}给 \(\mu*\mathbf1=\varepsilon\)。因此从 \(b=\mathbf1*a\) 卷积 \(\mu\) 得 \(a=\mu*b\)；反向卷积 \(\mathbf1\) 同样成立。
\end{proof}
此结果称为\term[mobius-fanyan]{Möbius 反演}[Möbius inversion]。例如 \(\varphi=\mu*\operatorname{id}\)，其中 \(\operatorname{id}(n)=n\)；由\eqref{b3:eq:25-log-divisors}还得 \(\Lambda=\mu*\log\)。这些恒等式是有限和的精确关系，后面乘进 \(n^{-s}\) 并无限求和时才出现收敛问题。

\ProseParagraphBreak

以 \(n=12\) 为例，只有平方自由约数 \(1,2,3,6\) 对反演公式有贡献。因此
\[
 \varphi(12)=12-6-4+2=4,
 \qquad
 \Lambda(12)=\log12-\log6-\log4+\log2=0.
\]
前式把所有与十二有公共素因子的整数逐次剔除，后式则使含两个不同素因子的对数贡献恰好抵消。对 \(n=p^r\)，后一式只剩 \(\log p^r-\log p^{r-1}=\log p\)。这些小例子显示反演既可用于计数，也可用来制造只集中在素数幂上的权重。

\begin{proposition}[平方自由筛选恒等式]
\label{b3:prop:25-squarefree-filter}
对每个正整数 \(n\)，
\[
 \mu(n)^2=\sum_{d^2\mid n}\mu(d).
\]
\end{proposition}
\begin{proof}
右侧的非零项对应那些平方仍整除 \(n\) 的不同素数的子集。如果共有 \(r\) 个这样的素数，和为 \((1-1)^r\)。当 \(n\) 平方自由时 \(r=0\)，值为一；否则值为零。这恰好等于左侧。
\end{proof}

这里的条件是 \(d^2\mid n\)，不是通常卷积中的 \(d\mid n\)。平方自由性要求排除素数平方，因此反演中的筛选条件也相应变成平方整除。把这项有限恒等式在 \(n\le x\) 上求和，随后交换两个有限和，便能把平方自由整数的计数转成 \(\lfloor x/d^2\rfloor\) 的估计；最后一节将算出其密度。

\ProseParagraphBreak

另一个有用的有限模型是多重因子计数。令 \(d_k(n)\) 表示有序正整数 \(k\) 元组 \((n_1,\ldots,n_k)\) 中满足 \(n_1\cdots n_k=n\) 的个数，则
\[
 d_k=\underbrace{\mathbf1*\cdots*\mathbf1}_{k\text{ 个因子}},
 \qquad d_k(p^r)=\binom{r+k-1}{k-1}.
\]
第二个公式来自把素数指数 \(r\) 分配到 \(k\) 个非负整数：在 \(r\) 个同类记号之间插入 \(k-1\) 个分隔符，共有所列组合数种分配。不同素数的指数分配独立，故 \(d_k\) 可乘。\(k=2\) 时恢复约数个数 \(r+1\)，而 \(k=1\) 时应理解为恒等于一的函数。
